\documentclass[11pt]{article}
\usepackage{amsmath,amssymb}
\usepackage[margin=1in]{geometry}

\title{Step 277: Closed-Form Search for \(\Phi_{\max}\)}
\author{Six Birds Foundations III}
\date{}

\begin{document}
\maketitle

\section*{Target}
The inherited Branch A value is
\[
\Phi_{\max}(0.35,2.0)\approx 0.4904766200298252.
\]
Step 277 tests whether this number has a small integer-relation expression in
standard constants.

\section*{Search}
The PSLQ basis included \(\pi\), powers and reciprocals of \(\pi\), \(e\),
logarithms \(\log 2,\log 3,\log\pi,\log(2\pi)\), Euler's constant,
\(\zeta(2),\zeta(3)\), Catalan's constant, Khinchin's constant,
\(\Gamma(1/4)\), \(\zeta(1/2)\), \(\zeta'(0)\), and Bessel-zero constants.

Transforms tested included
\[
\Phi,\quad \Phi^2,\quad 1-\Phi,\quad \log\Phi,\quad
\Phi\sqrt\pi,\quad \Phi/\log 2.
\]
Strict PSLQ used tolerance \(10^{-50}\) and coefficient bound \(10^4\).

\section*{Result}
No nontrivial integer relation involving \(\Phi\) was found.  Exact relations
with zero coefficient on \(\Phi\), such as \(\zeta(2)/\pi^2=1/6\), were rejected
as basis-only identities.

The closest simple hand candidate was
\[
\frac{\log 2}{\sqrt2},
\]
but its residual is approximately \(3.48\times10^{-4}\), so it fails.

\section*{Verdict}
\[
\boxed{\texttt{V\_branch\_A\_no\_closed\_form}}
\]
The value remains a provisional Sonine-family numerical constant within the
tested basis.  This is not a Branch A closure theorem and has no RH implication.

\end{document}
